

\addcontentsline{toc}{part}{Part I. The conjecture as a projector problem}

\chapter{The rational Hodge target}
\section{Statement and status}
Let $X$ be a smooth projective complex variety. For each codimension $p$, define
\[
\Hdg^p(X)=H^{2p}(X,\Q)\cap H^{p,p}(X)
\]
and
\[
\Alg^p(X)=\im\bigl(\CH^p(X)_{\Q}\xrightarrow{\cl}H^{2p}(X,\Q)\bigr).
\]
Every algebraic cycle class has type $(p,p)$, so $\Alg^p(X)\subseteq\Hdg^p(X)$. The rational Hodge conjecture asserts equality for every smooth projective $X$ and every $p$ \cite{ClayHodge,VoisinHodge}.

Choose a positive inner product on $H^{2p}(X,\R)$. The orthogonal projectors $P_{\Hdg}^p$ and $P_{\Alg}^p$ act on the real scalar extensions $\Hdg^p(X)\otimes_{\Q}\R$ and $\Alg^p(X)\otimes_{\Q}\R$. These real orthogonal projectors provide a reformulation. They are not asserted to be rational matrices or algebraic correspondences. Since the algebraic space is nested in the Hodge space,
\[
P_{\Hdg}^pP_{\Alg}^p=P_{\Alg}^p.
\]

\begin{proposition}[Master obstruction projector]
The operator
\[
\mathcal O_X^p=P_{\Hdg}^p-P_{\Alg}^p
\]
is the orthogonal projector onto $(\Hdg^p(X)\otimes_{\Q}\R)\cap(\Alg^p(X)\otimes_{\Q}\R)^\perp$. Therefore
\[
\mathcal O_X^p=0
\quad\Longleftrightarrow\quad
\Hdg^p(X)=\Alg^p(X).
\]
\end{proposition}
\begin{proof}
For nested orthogonal projectors $Q\le P$, the difference $P-Q$ is self-adjoint and
\[
(P-Q)^2=P-PQ-QP+Q=P-Q.
\]
Its range is $\im P\cap(\im Q)^\perp$.
\end{proof}

\begin{figure}[tbp]\centering\includegraphics[width=0.98\linewidth]{04_projector_diagnostic.pdf}\caption{The real orthogonal obstruction is a diagnostic reformulation. An algebraic realization of the full Hodge projector is not supplied.}\end{figure}

\section{Why dimension four is the first focused target}
The first general unresolved middle case occurs for primitive degree-four cohomology on a smooth projective fourfold. The Hodge diamond highlights the central packet.

\begin{figure}[tbp]\centering\includegraphics[width=0.98\linewidth]{03_fourfold_middle.pdf}\caption{Hodge bidegree positions for a fourfold. The diagram gives indices rather than numerical Hodge dimensions.}\end{figure}

\chapter{Primitive and coniveau reductions}
\section{Lefschetz packet decomposition}
Let $h\in H^2(X,\Q)$ be the class of an ample divisor and let $L(\alpha)=h\smile\alpha$. Hard Lefschetz and primitive decomposition yield
\[
H^{2p}(X,\Q)=\bigoplus_{r\ge0}L^rH_{\mathrm{prim}}^{2p-2r}(X,\Q).
\]
The decomposition respects Hodge type.

\begin{theorem}[Primitive reduction]
If every primitive rational Hodge class is algebraic, then every rational Hodge class is algebraic.
\end{theorem}
\begin{proof}
Write $\alpha=\sum_rL^r\alpha_r$ with each $\alpha_r$ primitive and of Hodge type. The class $h$ is algebraic. If $\alpha_r$ is algebraic, then $L^r\alpha_r=h^r\smile\alpha_r$ is algebraic. Summing gives an algebraic representative of $\alpha$.
\end{proof}

For a fourfold and $p=2$,
\[
\alpha=L\beta+\alpha_{\mathrm{prim}},
\qquad
\beta\in H^{1,1}(X)\cap H^2(X,\Q).
\]
The Lefschetz $(1,1)$ theorem makes $\beta$ algebraic. Only $\alpha_{\mathrm{prim}}$ remains.



\section{Support localization and coniveau}
Let $N^pH^{2p}(X,\Q)$ be the geometric coniveau-$p$ subspace. In degree $2p$, purity forces support contributions of codimension $p$ to be generated by fundamental classes of codimension-$p$ components.

\begin{theorem}[Degree-$2p$ coniveau identification]
For smooth projective $X$,
\[
N^pH^{2p}(X,\Q)=\Alg^p(X).
\]
\end{theorem}
\begin{proof}[Proof sketch]
The inclusion from right to left is immediate because a codimension-$p$ cycle is supported on its underlying algebraic set. Conversely, a class in coniveau $p$ is in the image of cohomology with supports in an algebraic subset $Z$ of codimension at least $p$. Contributions from components of codimension greater than $p$ vanish in degree $2p$. Purity identifies the degree-$2p$ supported contribution of each codimension-$p$ irreducible component with its fundamental class. Resolution and additivity reduce a general $Z$ to these components.
\end{proof}

Thus the conjecture is equivalently the support statement
\[
\Hdg^p(X)\subseteq N^pH^{2p}(X,\Q).
\]



\addcontentsline{toc}{part}{Part II. Algebraic cycle frames and exact polyphase refinement}

\chapter{Positive cycle-frame calculus}
\section{Analysis and synthesis}
Fix a finite family of rational algebraic cycle classes
\[
c_1,\ldots,c_N\in H^{2p}(X,\Q).
\]
Let $V\subseteq H^{2p}(X,\R)$ be the real extension of the rational Hodge packet under study. For the positive frame certificate below, require every $c_j$ to belong to $V$. This can be achieved by a rational algebraic packet projector when such a projector has been constructed. Projection through an arbitrary orthogonal Hodge projector does not establish that the projected vectors are algebraic. Define
\[
C:\R^N\to H^{2p}(X,\R),
\qquad Ce_j=c_j,
\]
and the positive frame operator
\[
S=CC^*=\sum_{j=1}^N|c_j\rangle\langle c_j|.
\]

\begin{theorem}[Finite positive cycle-frame certificate]
Let $A=\Span_{\R}\{c_1,\ldots,c_N\}$. Then
\[
\im S=A,
\qquad
\ker S=A^\perp.
\]
Under the stated containment $c_j\in V$, the selected cycles span $V$ if and only if $S|_V$ is positive definite. Without containment, positivity only detects the absence of vectors in $V$ orthogonal to the cycle span and does not prove that span contains $V$.
\end{theorem}
\begin{proof}
The equality $\im(CC^*)=\im C$ is standard. Moreover
\[
\langle v,Sv\rangle=\|C^*v\|^2=\sum_j|\langle c_j,v\rangle|^2.
\]
This vanishes exactly when $v$ is orthogonal to every cycle class.
\end{proof}

If $H$ is a matrix whose columns form an orthonormal basis of $V$, the finite Gram certificate is
\[
G=H^*SH.
\]
Under $A\subseteq V$, we have $A=V$ if and only if $\lambda_{\min}(G)>0$. Without that containment, use a rational basis matrix $H_{\Q}$ for the target and the exact certificate $\operatorname{rank}_{\Q}[C\ H_{\Q}]=\operatorname{rank}_{\Q}C$ to test $V\subseteq A$.

\begin{figure}[tbp]\centering\includegraphics[width=0.98\linewidth]{06_corrected_rational_frame.pdf}\caption{A finite rational frame example with its cycle vectors contained in the selected target. Actual geometric applications must supply algebraic vectors with verified cohomology coordinates.}\end{figure}



The orthogonal synthesis of a class $\alpha$ from the cycle span is
\[
P_A\alpha=S S^\dagger\alpha.
\]
If $\alpha\in A$, this is exact.



\section{Packet assembly without false orthogonality}
Partition the cycle indices into packets $I_\lambda$. Define
\[
S_\lambda=\sum_{j\in I_\lambda}|c_j\rangle\langle c_j|.
\]
Then
\[
S=\sum_\lambda S_\lambda.
\]
No orthogonality of the child spans is assumed.

\begin{theorem}[Positive child-kernel intersection]
For any subspace $V$,
\[
\ker(S|_V)=\bigcap_\lambda\ker(S_\lambda|_V).
\]
\end{theorem}
\begin{proof}
If $Sv=0$, positivity gives
\[
0=\langle v,Sv\rangle=\sum_\lambda\langle v,S_\lambda v\rangle.
\]
Each nonnegative summand vanishes, which implies $S_\lambda v=0$. The converse is immediate.
\end{proof}

\begin{figure}[tbp]\centering\includegraphics[width=0.98\linewidth]{07_positive_polyphase.pdf}\caption{Positive parent-child frame assembly. The identity combines frames without assuming their ranges are orthogonal.}\end{figure}

\chapter{Degree-residue polyphase refinement}
\section{Exact parent-child frame law}
Fix a projective embedding and assign each cycle its degree. For modulus $M$, define the residue packet
\[
I_{M,r}=\{j:\deg Z_j\equiv r\pmod M\},
\qquad
S_{M,r}=\sum_{j\in I_{M,r}}|c_j\rangle\langle c_j|.
\]

\begin{theorem}[Exact degree polyphase refinement]
For every positive integer $L$,
\[
S_{M,r}=\sum_{\ell=0}^{L-1}S_{LM,r+\ell M}.
\]
\end{theorem}
\begin{proof}
The congruence class $r$ modulo $M$ is the disjoint union of the classes $r+\ell M$ modulo $LM$. The rank-one frame terms are partitioned accordingly.
\end{proof}





\section{Multiaxial packet addresses}
Degree residue is only one address. A refined cycle packet may be labeled by
\[
\lambda=(r_{\mathrm{Lef}},r_{\deg},\chi_{\mathrm{mon}},s_{\mathrm{supp}},\nu_{\mathrm{Frob}}).
\]
The entries record Lefschetz depth, degree residue, monodromy or symmetry type, support stratum, and arithmetic phase. The total positive frame remains the sum of all leaves. This allows independent methods to solve different geometric sectors before exact synthesis.

\chapter{Finite saturation and exact rational closure}
For each degree cutoff $D$, let $A_D^p(X)$ be the rational span of codimension-$p$ cycle classes of degree at most $D$. Then
\[
A_1^p\subseteq A_2^p\subseteq\cdots\subseteq H^{2p}(X,\Q).
\]

\begin{theorem}[Finite stabilization]
There exists $D_0$ such that $A_D^p(X)=A_{D_0}^p(X)$ for every $D\ge D_0$.
\end{theorem}
\begin{proof}
The dimensions form a nondecreasing sequence of integers bounded by $\dim H^{2p}(X,\Q)$. Only finitely many strict increases are possible.
\end{proof}

\begin{corollary}[Approximation becomes exact]
Let $P_D$ be the orthogonal projector onto $A_D^p(X)\otimes\R$. For a rational class $\alpha$,
\[
\lim_{D\to\infty}\|(I-P_D)\alpha\|=0
\]
if and only if $\alpha\in A_D^p(X)$ for some finite $D$.
\end{corollary}
\begin{proof}
After stabilization, the residual is constant. A constant tail converging to zero is zero.
\end{proof}

Because the matrices and target vector can be chosen over $\Q$, real membership implies rational membership by row reduction over $\Q$.



\begin{resultbox}[Finite completion certificate]
If $H$ is a rational basis matrix of $\Hdg^p(X)$ and $C_D$ is the cycle matrix, then the Hodge conjecture for $(X,p)$ is equivalent to the existence of $D$ with
\[
\rank_{\Q} C_D=\rank_{\Q}[C_D\ H].
\]
A positive compressed frame is equivalent only when the supplied cycle vectors lie in the target packet, or when a separately verified invariant restriction supplies the required containment. The rank equality is valid without that extra hypothesis.
\end{resultbox}

\addcontentsline{toc}{part}{Part III. Support packets and correspondence reconstruction}

\chapter{Support verification and multiplicity recovery}
Let $Z\subseteq X$ be a closed algebraic subset and $j:X\setminus Z\hookrightarrow X$. The support exact sequence contains
\[
H_Z^{2p}(X,\Q)\longrightarrow H^{2p}(X,\Q)\xrightarrow{j^*}H^{2p}(X\setminus Z,\Q).
\]

\begin{theorem}[Exact support test]
A class $\alpha\in H^{2p}(X,\Q)$ is supported on $Z$ if and only if $j^*\alpha=0$.
\end{theorem}
\begin{proof}
This is exactness at the middle term.
\end{proof}

\begin{figure}[tbp]\centering\includegraphics[width=0.98\linewidth]{08_support_verification.pdf}\caption{The restriction-to-complement test verifies a supplied support. It does not construct a support for an arbitrary Hodge class.}\end{figure}

Suppose $Z$ has codimension at least $p$ and its codimension-$p$ components are $Z_1,\ldots,Z_r$. Purity gives
\[
\alpha=\sum_{i=1}^r q_i[Z_i].
\]
Choose complementary test classes $\gamma_k$ and define
\[
M_{ki}=\int_X[Z_i]\smile\gamma_k,
\qquad
b_k=\int_X\alpha\smile\gamma_k.
\]
Then $Mq=b$. Removing dependent columns yields exact rational multiplicities.



\chapter{Divisor-lowering correspondence channels}
\section{Degree bookkeeping}
Let $X$ have complex dimension $n$, let $V\subseteq H_{\mathrm{prim}}^{2p}(X,\Q)$ be a rational Hodge substructure, and let $Y_r$ be smooth projective. Choose correspondences
\[
\Gamma_r\in\CH^{n-p+1}(X\times Y_r)_{\Q},
\qquad
\Lambda_r\in\CH^{\dim Y_r+p-1}(Y_r\times X)_{\Q}.
\]
They induce
\[
A_r=(\Gamma_r)_*:V\to H^2(Y_r,\Q),
\qquad
B_r=(\Lambda_r)_*:H^2(Y_r,\Q)\to H^{2p}(X,\Q).
\]
Algebraic correspondences preserve rationality and shift Hodge type by the required bidegree.

\begin{theorem}[Invertible divisor-lowering transfer]
Suppose
\[
T=\sum_{r=1}^M B_rA_r
\]
preserves $V$ and is invertible on $V$. Then every rational Hodge class in $V$ is algebraic.
\end{theorem}
\begin{proof}
Let $\alpha\in V$ be of type $(p,p)$. Each $A_r\alpha$ is rational of type $(1,1)$ on $Y_r$, hence is the class of a rational divisor $D_r$ by the Lefschetz $(1,1)$ theorem. Therefore
\[
T\alpha=\sum_r(\Lambda_r)_*[D_r]
\]
is algebraic on $X$.

Let $\chi_T(t)=t^d+c_{d-1}t^{d-1}+\cdots+c_0$ be the characteristic polynomial on $V$. Invertibility means $c_0\ne0$. Cayley-Hamilton gives
\[
T^{-1}=-\frac{1}{c_0}\bigl(T^{d-1}+c_{d-1}T^{d-2}+\cdots+c_1I\bigr).
\]
Each power of $T$ is induced by compositions and rational sums of algebraic correspondences. Thus $T^{-1}$ preserves algebraic cycle classes. Hence $\alpha=T^{-1}(T\alpha)$ is algebraic.
\end{proof}

\begin{figure}[tbp]\centering\includegraphics[width=0.98\linewidth]{09_correspondence_degrees.pdf}\caption{The exact degree shifts required for lowering to divisor degree and returning. Actual Chow kernels and invertibility are additional inputs.}\end{figure}

\begin{figure}[tbp]\centering\includegraphics[width=0.98\linewidth]{10_rational_inverse.pdf}\caption{A finite rational example of polynomial inversion. The geometric theorem requires the operator to be induced by algebraic correspondences.}\end{figure}

\section{Positive multi-channel form}
Assume the polarized adjoint $A_r^\dagger$ is realized by an algebraic return correspondence $B_r$. This is an additional geometric hypothesis, not an automatic consequence in arbitrary settings. Then
\[
T=\sum_rA_r^\dagger A_r
\]
is positive and
\[
\langle v,Tv\rangle=\sum_r\|A_rv\|^2.
\]
Therefore
\[
\ker T=\bigcap_r\ker A_r.
\]
Joint injectivity of the lowering channels is enough.





\begin{corollary}[Finite Gram certificate]
Choose a rational basis $e_1,\ldots,e_d$ of $V$. If
\[
G_{ij}=\sum_r\langle A_re_i,A_re_j\rangle_{Y_r}
\]
has nonzero determinant, every rational Hodge class in $V$ is algebraic.
\end{corollary}

\chapter{Family propagation}
Let $\pi:\mathcal X\to S$ and $\rho_r:\mathcal Y_r\to S$ be smooth projective families over a connected base. Relative algebraic correspondences induce morphisms of Gauss-Manin local systems. In a flat rational trivialization, their transfer matrix is constant.

\begin{theorem}[Propagation of invertibility]
If a relative transfer $T_s$ is invertible on one fiber $s_0$, it is invertible on every fiber in the connected base on which the local systems and correspondences are defined.
\end{theorem}
\begin{proof}
The determinant of the constant rational matrix representing the local-system morphism is locally constant. It is nonzero at $s_0$, hence everywhere on the connected base.
\end{proof}



\addcontentsline{toc}{part}{Part IV. Closed branches and classical validation cases}

\chapter{Geometric closure rules}
\section{Affine pavings}
\begin{theorem}[Affine-paving branch]
If $X$ admits a finite filtration whose successive differences are affine spaces, then the cycle class map
\[
\CH^p(X)_{\Q}\to H^{2p}(X,\Q)
\]
is surjective for every $p$.
\end{theorem}
\begin{proof}[Proof sketch]
Induct through the filtration using localization for Chow groups and cohomology. Each affine cell contributes one even-degree generator represented by the closure of that cell. No odd cohomology is introduced.
\end{proof}

\section{Projective bundles, blow-ups, and descent}
\begin{proposition}[Projective-bundle closure]
If the rational Hodge conjecture holds for $X$, it holds for $\mathbb P(E)\to X$.
\end{proposition}
\begin{proof}
The projective-bundle formula expresses cohomology as a direct sum of powers of the algebraic hyperplane class times pullbacks of cohomology from $X$.
\end{proof}

\begin{proposition}[Blow-up closure]
If it holds for $X$ and a smooth center $Z$, it holds for $\operatorname{Bl}_ZX$.
\end{proposition}
\begin{proof}
Use the standard blow-up cohomology decomposition. The new summands are Tate twists of cohomology from $Z$ multiplied by powers of the algebraic exceptional divisor.
\end{proof}

\begin{proposition}[Generically finite descent]
Let $f:Y\to X$ be proper and generically finite of degree $d$ between smooth projective varieties of equal dimension. If the conjecture holds for $Y$, it holds for $X$.
\end{proposition}
\begin{proof}
For a Hodge class $\alpha$ on $X$, the pullback $f^*\alpha$ is algebraic on $Y$. Push forward and use $f_*f^*\alpha=d\alpha$.
\end{proof}



\chapter{Cubic fourfolds as a validation of the transfer theorem}
Let $X\subset\mathbb P^5$ be a smooth cubic fourfold and let $F(X)$ be its Fano variety of lines. The universal incidence variety
\[
P=\{(\ell,x):x\in\ell\}\subset F(X)\times X
\]
induces the Abel-Jacobi map
\[
A=p_*q^*:H^4_{\mathrm{prim}}(X,\Q)\to H^2_{\mathrm{prim}}(F(X),\Q).
\]
Beauville and Donagi proved that this is a polarized Hodge isomorphism up to the standard twist \cite{BeauvilleDonagi}. The polarized adjoint is represented by an algebraic incidence return after inserting the appropriate power of the Pluecker class. Therefore the divisor-lowering transfer is invertible.

\begin{resultbox}[Classical validation]
A primitive rational $(2,2)$ class on a cubic fourfold is sent to a rational $(1,1)$ class on $F(X)$, represented by a divisor. The algebraic return correspondence and its polynomial inverse reconstruct an algebraic codimension-two class on $X$. This is a validation of the packet criterion, not a claim of a new proof of the classical cubic-fourfold result.
\end{resultbox}





\chapter{Gushel-Mukai validation channels}
Debarre and Kuznetsov construct line and plane incidence correspondences relating vanishing middle cohomology of suitable Gushel-Mukai varieties to primitive degree-two cohomology of associated double EPW sextics \cite{DebarreKuznetsov}. In the ranges where the polarized Hodge isomorphisms are established, the same divisor-lowering theorem applies.

\begin{itemize}
\item For suitable GM fourfolds, the vanishing $H^4$ packet is lowered by universal lines to divisor degree.
\item For suitable GM sixfolds, the vanishing $H^6$ packet is lowered by universal planes to divisor degree.
\item Ambient and tautological classes are treated separately and are algebraic.
\end{itemize}



\addcontentsline{toc}{part}{Part V. The analytic and K-theoretic branch}

\chapter{Hodge packets in topological K-theory}
For finite CW complexes, the rational topological Chern character gives
\[
\operatorname{ch}_{\mathrm{top}}:K^0_{\mathrm{top}}(X)\otimes\Q\xrightarrow{\sim}H^{\mathrm{even}}(X,\Q).
\]
For smooth projective $X$, the algebraic Chern character identifies algebraic $K$-theory rationally with the Chow ring after the usual filtration. Thus the rational Hodge conjecture for all degrees is equivalent to algebraization of topological $K$-classes whose Chern-character components are rational Hodge classes \cite{AtiyahHirzebruch}.

\section{Adams degree projectors}
For integer $m\ge2$, Adams operations satisfy
\[
\operatorname{ch}_q(\psi^m\xi)=m^q\operatorname{ch}_q(\xi).
\]
Define
\[
e_p^{(m)}=\prod_{\substack{0\le q\le\dim X\\q\ne p}}
\frac{\psi^m-m^qI}{m^p-m^q}.
\]
Then
\[
\operatorname{ch}_q(e_p^{(m)}\xi)=\delta_{pq}\operatorname{ch}_p(\xi).
\]
This exactly isolates codimension $p$ inside a topological $K$-packet.



\chapter{Nonlinear Dolbeault obstruction packets}
Let $E$ be a smooth complex bundle and consider a candidate $(0,1)$-connection. After gauge fixing, write the integrability equation abstractly as
\[
F(a)=f+La+B(a,a)=0,
\]
where $L$ is the linearized operator and $B$ is a bounded bilinear term. Assume a Fredholm splitting of the target
\[
I=Q+\Pi,
\]
where $Q$ projects onto the closed range of $L$ and $\Pi$ onto a finite-dimensional obstruction space. Choose a gauge-fixed complement to $\ker L$ as the domain and let $R$ be a right inverse of $L$ taking values in that complement. Use norms on the projected data for which $\|QB(a,b)\|\le C\|a\|\|b\|$.

\begin{theorem}[Range-packet contraction]
Assume
\[
\|R\|\le M,
\qquad
\|B(a,b)\|\le C\|a\|\|b\|,
\qquad
4CM^2\|Qf\|<1.
\]
Then the map
\[
\mathcal T(a)=-RQ(f+B(a,a))
\]
has a unique fixed point in the gauge-fixed ball $\|a\|\le2M\|Qf\|$. If $Qf=0$, the ball is the singleton $\{0\}$ and the range solution is immediate. At the fixed point, the range component of $F(a)$ vanishes. The remaining obstruction is
\[
\Phi=\Pi(f+B(a,a)).
\]
The full integrability equation holds if and only if $\Phi=0$.
\end{theorem}
\begin{proof}
For $\|a\|\le r=2M\|Qf\|$,
\[
\|\mathcal T(a)\|\le M\|Qf\|+MC r^2\le r
\]
under the stated bound. Moreover
\[
\|\mathcal T(a)-\mathcal T(b)\|
\le MC(\|a\|+\|b\|)\|a-b\|.
\]
The Lipschitz constant is strictly less than one on the stated ball.
Banach's fixed-point theorem gives the solution. Applying $Q$ to the equation proves range cancellation. The remaining component is precisely $\Phi$.
\end{proof}





\section{A positive spectral microscope}
Large ample twists can make integrability defects visible in the asymptotic degree distribution of Dolbeault-Dirac kernels. Charbonneau and Stern obtain asymptotic Hodge filtrations in which positive coefficients depend on powers of $F_A^{0,2}$ \cite{CharbonneauStern}. Schematically,
\[
\Tr(P_{2j}\Pi_{A,k})
=c_{n,j}k^{n-2j}\|(F_A^{0,2})^{\wedge j}\|_{L^2}^2+O(k^{n-2j-1}),
\qquad c_{n,j}>0.
\]
This detects nonintegrability. It does not by itself remove the defect.



\section{Trace blindness and the true missing bridge}
Chern characters control traces of powers of curvature. An endomorphism-valued obstruction can be nonzero while all such trace moments vanish. For example, $N=\left(\begin{smallmatrix}0&1&0\\0&0&1\\0&0&0\end{smallmatrix}\right)$ is nonzero, while $N^3=0$ and $\operatorname{Tr}(N^r)=0$ for every positive integer $r$. This is the elementary negative control.



\begin{openproblem}[Analytic bridge]
The unresolved analytic theorem is not the range equation. It is a positive saturation statement forcing the finite obstruction packet $\Phi$ to vanish for every rational Hodge $K$-packet after allowable stabilization. Proving that statement in general would be essentially a K-theoretic form of the Hodge conjecture.
\end{openproblem}

\section{Consequences of the obstruction split}
The obstruction split prevents a common logical error. Solving the projected range equation does not solve the original integrability problem unless the complementary packet also vanishes. The calculus therefore records three distinct outcomes. First, if $\Pi=0$, the problem is unobstructed and the contraction theorem gives an actual integrable structure. Second, if $\Pi\ne0$ but $\Phi=0$ for geometric reasons, the correction again closes. Third, if $\Phi\ne0$, the method has produced a finite certificate explaining exactly why that representative cannot be algebraized by the chosen deformation.

This creates a practical research program. One can refine the finite obstruction space by symmetry, monodromy, support, or arithmetic phase, obtaining
\[
\Phi=\sum_{\lambda}\Phi_{\lambda}.
\]
Since the decomposition is complete, integrability is equivalent to the simultaneous vanishing of all children. Different geometric arguments may eliminate different children. A symmetry packet may vanish by representation theory, a support packet by localization, and an arithmetic packet by compatibility across reductions. The analytic branch therefore remains useful even before the general bridge is proved: it supplies exact special-case theorems whenever all obstruction children can be eliminated.

The negative controls are equally important. Ample twisting can amplify $F_A^{0,2}$ in a spectral observable without reducing it. Direct-sum stabilization adds positive defect energies for block-diagonal connections. Trace identities may miss nilpotent components. These facts rule out several tempting but invalid shortcuts and keep the remaining theorem sharply stated.


\addcontentsline{toc}{part}{Part VI. Atomic status ledger and focused next theorem}

\chapter{Dependency graph and proof status}
The project separates proved statements, classical inputs, conditional mechanisms, and open bridges.



\begin{longtable}{p{0.40\textwidth}p{0.16\textwidth}p{0.34\textwidth}}
\caption{Atomic status ledger.}\label{tab:status}\\
\toprule
Statement & Status & Boundary \\
\midrule
\endfirsthead
\toprule
Statement & Status & Boundary \\
\midrule
\endhead
Master obstruction projector & Proved & Finite-dimensional linear algebra \\
Primitive Lefschetz reduction & Proved & Uses hard Lefschetz and algebraicity of the polarization \\
Degree-$2p$ coniveau identification & Proved & Rational coefficients and purity \\
Positive cycle-frame certificate & Corrected & Requires explicit algebraic cycles contained in the selected packet \\
Exact degree polyphase assembly & Proved & Frame operators, not sums of span projectors \\
Finite saturation of bounded-degree spans & Proved & Non-effective degree bound \\
Support verification & Proved & Candidate support must already be supplied \\
Multiplicity reconstruction & Proved & Finite rational linear algebra \\
Divisor-lowering transfer theorem & Proved & Requires explicit algebraic correspondences and invertibility \\
Positive multi-channel criterion & Conditional & Algebraic realization of polarized adjoints is an extra hypothesis \\
Family propagation & Proved & Relative correspondences over a connected smooth base \\
Affine-paving branch & Proved & Classical localization argument \\
Projective bundle, blow-up, descent rules & Proved & Standard cohomological decompositions \\
Cubic and GM channels & Classical validation & Depend on established incidence Hodge isomorphisms \\
Adams degree projector & Proved & Rational topological K-theory \\
Range-packet nonlinear correction & Proved under hypotheses & Closed range, right inverse, and contraction bound \\
Vanishing of every harmonic obstruction packet & Open & Main analytic bridge \\
Primitive fourfold multi-incidence saturation & Open & Main geometric target \\
General rational Hodge conjecture & Open & No general cycle-generation theorem is claimed \\
\bottomrule
\end{longtable}



\chapter{A focused multi-incidence target}
Let $X$ be a special smooth projective fourfold. Construct algebraic lowering channels from families of subvarieties:
\[
A_{\mathrm{line}},
A_{\mathrm{conic}},
A_{\mathrm{plane}},
A_{\mathrm{exceptional}}.
\]
The exact target is
\[
\ker A_{\mathrm{line}}
\cap\ker A_{\mathrm{conic}}
\cap\ker A_{\mathrm{plane}}
\cap\ker A_{\mathrm{exceptional}}=\{0\}
\]
on primitive degree-four cohomology. Under algebraic adjoint hypotheses, equivalently prove
\[
T_{\mathrm{tot}}
=\sum_rA_r^\dagger A_r
\quad\text{satisfies}\quad
\lambda_{\min}(T_{\mathrm{tot}})>0.
\]



\section{Why this target is realistic}
The problem is finite after the geometry is chosen:
\begin{enumerate}
\item compute the primitive cohomology representation and its symmetry decomposition
\item construct the incidence correspondences and their action
\item split the action into monodromy or symmetry packets
\item compute the joint kernel or positive Gram matrix over $\Q$
\item prove exact nonvanishing through rational rank, representation theory, or a dual certificate
\item invoke the divisor-lowering theorem to reconstruct cycles
\end{enumerate}
No single incidence family must be complete. Polyphase analysis permits separate geometric mechanisms to occupy separate primitive sectors.


\section{Three exact certificate forms}
Once the incidence channels have been constructed, the saturation statement admits three finite certificates. The geometric form is the kernel identity
\[
\bigcap_r\ker A_r=\{0\}.
\]
When algebraic correspondences realize the polarized adjoints, the equivalent positive form is
\[
\lambda_{\min}\left(\sum_rA_r^\dagger A_r\right)>0.
\]
In a rational basis, the final exact proof object is full column rank of the vertically stacked channel matrix. The kernel form explains the geometry, the positive form is stable for numerical scouting, and the rational rank form is appropriate for symbolic certification.

Two negative controls should accompany the proof. First, remove one channel at a time and recompute the joint kernel. If the kernel grows, the omitted incidence family is genuinely necessary. Second, refine a parent channel into geometric children and verify that their positive transfers sum exactly to the parent. This separates true packet coverage from accidental numerical invertibility.

A complete result should identify which primitive symmetry packets are detected by lines, conics, planes, or exceptional strata. The proof then becomes a finite occupancy statement: every primitive packet must be detected at least once. That structural description is more informative than a determinant alone and is the form most likely to propagate through families and degenerations.

\chapter{Computational and formal proof workflow}
Numerical experiments can scout candidate families but cannot replace an exact rank certificate. The intended pipeline is:
\[
\begin{aligned}
\text{exact incidence data}
&\longrightarrow \text{rational cohomology basis}
\longrightarrow \text{packet transfer matrices}\\
&\longrightarrow \text{rank or positivity certificate}
\longrightarrow \text{symbolic theorem proof}.
\end{aligned}
\]



\section{Reproducible finite checks}
The accompanying Python generator produces all figures and runs eighteen independent finite checks covering:
\begin{itemize}
\item frame idempotency and range rank
\item positive Hodge Gram spectrum
\item exact modulo-$2$ and modulo-$4$ polyphase assembly
\item positive transfer rank and energy identities
\item Cayley-Hamilton inversion
\item contraction of the removable nonlinear packet
\item Adams projector idempotency and degree isolation
\item nilpotent trace blindness
\item support multiplicity recovery
\item finite rank stabilization
\end{itemize}
These tests verify the finite examples and implementation. They do not establish the open geometric saturation theorem.

\compactpart{Part VII. Fermat Fourfold Character Packets and the Constructive Program}

\chapter{Correct scope and claim boundary}
The projector-resolved Hodge calculus naturally specializes to Fermat fourfolds
\[
X_m^4=\{x_0^m+x_1^m+x_2^m+x_3^m+x_4^m+x_5^m=0\}\subset\mathbf P^5,
\qquad m\ge 3.
\]
The diagonal symmetry group makes the primitive middle cohomology unusually explicit: it splits into one-dimensional complex character lines, and rational Hodge packets are finite Galois orbits of those lines. This permits exact symbolic analysis of individual addresses, conductors, local prime phases, and candidate algebraic correspondences. A complete classification at every mixed conductor is a separate coverage problem.

\begin{claimboundary}
The general rational Hodge conjecture remains open. For Fermat fourfolds, however, Kang is cited in the checked later literature as proving the relevant generalized Hodge statement in arbitrary degree \cite{KangFermat}. This complete edition therefore distinguishes two questions:
\begin{enumerate}
\item \emph{Existence}: does every rational Fermat $(2,2)$ class have an algebraic representative?
\item \emph{Construction}: can one write an explicit packet-resolved algebraic cycle or correspondence for every character orbit?
\end{enumerate}
The present Fermat program is primarily constructive. Its new proof objects are conductor projectors, exact packet identities, explicit Delsarte maps, weighted dyadic refinements, the proposed all-degree trivalent-tree classification, CRT--Kummer phase factorizations, and explicit CM carrier data. The relation between that existence theorem and later packet-specific combinatorial work should be stated explicitly in any publication.
\end{claimboundary}



The constructive objective is
\begin{resultbox}[Packet-resolved construction target]
For every rational packet
\[
W\subset H^4_{\mathrm{prim}}(X_m^4,\Q)\cap H^{2,2},
\]
construct
\[
Z_W\in\CH^2(X_m^4)_\Q
\qquad\text{with}\qquad
[Z_W]\ne0\text{ in }W.
\]
\end{resultbox}
Once one nonzero algebraic vector is produced in a rational irreducible diagonal-symmetry packet, the entire rational orbit is algebraic.

\chapter{Universal balanced-packet equations}
Let
\[
U_m=(\mathbf Z/m\mathbf Z)^\times
\]
and let \(\langle r\rangle_m\in\{1,\ldots,m\}\) denote the least positive residue, with \(\langle r\rangle_m=m\) when \(m\mid r\). A primitive character address is a six-tuple
\[
a=(a_0,\ldots,a_5),\qquad 1\le a_i<m,
\qquad \sum_i a_i\equiv0\pmod m.
\]
The Hodge level of the conjugate indexed by \(u\in U_m\) is
\[
\ell_a(u)=\frac1m\sum_{i=0}^5\langle ua_i\rangle_m.
\]
The character line has type \((2,2)\) exactly when \(\ell_a(1)=3\). Its rational Galois orbit is a Hodge packet exactly when
\begin{equation}\label{eq:legacy-7.1}
\boxed{
\ell_a(u)=3\qquad\forall u\in U_m.
}\end{equation}

To remove coordinate permutations, introduce the count vector
\[
n(a)=(n_1,\ldots,n_{m-1})\in\mathbf N^{m-1},
\qquad n_r=\#\{i:a_i=r\}.
\]
Let
\[
R_m(u,r)=\langle ur\rangle_m.
\]
Then the complete balanced six-packet space is the finite integer set
\begin{equation}\label{eq:legacy-7.2}
\boxed{
\mathcal B_m=
\left\{
 n\in\mathbf N^{m-1}:
 \mathbf1^Tn=6,
 \quad R_mn=3m\,\mathbf1
\right\}.
}\end{equation}
The unit group permutes the residue coordinates. After additionally identifying coordinate permutations by algebraic automorphisms, the count-vector descriptions correspond to \(U_m\)-orbits in \(\mathcal B_m\). The dimension of a fixed rational character packet is computed from distinct ordered unit conjugates before this permutation quotient.



\section{Centered Fourier form}
Define
\[
\rho_m(r)=\langle r\rangle_m-\frac m2,
\qquad
F_a(u)=\sum_i\rho_m(ua_i).
\]
Then \(a\) is balanced if and only if \(F_a(u)=0\) for all units. Moreover,
\[
F_a(-u)=-F_a(u),
\]
so only odd characters occur. For a Dirichlet character \(\chi\) on \(U_m\), put
\[
\mathcal B_{m,\chi}(r)=
\sum_{u\in U_m}\rho_m(ur)\overline{\chi(u)}.
\]
Finite Fourier inversion gives the exact criterion
\begin{equation}\label{eq:legacy-7.3}
\boxed{
 n(a)\in\mathcal B_m
 \Longleftrightarrow
 \sum_i\mathcal B_{m,\chi}(a_i)=0
 \quad\text{for every odd }\chi.
}\end{equation}
Here $n(a)$ is the residue count vector. This formulation underlies the prime-power filtration and the proposed conductor-defect tree.

\begin{figure}[tbp]\centering\includegraphics[width=0.98\linewidth]{11_balanced_orbit70.pdf}\caption{Every ordered unit conjugate of the named degree-70 character has sum 210. This exact example identifies Hodge type and does not itself construct a cycle.}\end{figure}

\chapter{Exact conductor polyphase}
For a packet \(a\), define
\[
g(a)=\gcd(m,a_0,\ldots,a_5),
\qquad
c(a)=\frac{m}{g(a)}.
\]
For a count vector, $L_k$ denotes scaling each residue entry by $k$ while multiplying the modulus by $k$. Writing \(a=g(a)b\), the coordinate-power map
\[
\pi_{g(a)}:X_m^4\longrightarrow X_{c(a)}^4,
\qquad [x_i]\longmapsto[x_i^{g(a)}],
\]
sends the degree-\(c(a)\) character packet \(b\) to the degree-\(m\) packet \(a\).

\begin{theorem}[Conductor descent]
Every imprimitive balanced packet at degree \(m\) is the coordinate-power pullback of a unique full-conductor packet at its exact conductor \(c\mid m\). Consequently
\[
\mathcal B_m=\bigsqcup_{c\mid m}L_{m/c}\mathcal B_c^{\mathrm{full}}.
\]
\end{theorem}
\begin{proof}
If every occupied residue is divisible by \(g\), divide the six entries and the modulus by \(g\). Since
\[
\langle ugb_i\rangle_{gc}=g\langle \bar u b_i\rangle_c,
\]
the balance equation descends. The maximal common divisor is unique. The converse follows by scaling.
\end{proof}

For \(d\mid m\), let
\[
E_{m,d}=\frac1{(m/d)^5}\pi_{m/d}^*\pi_{m/d*}.
\]
On a character of exact conductor \(c\),
\[
E_{m,d}=\begin{cases}1,&c\mid d,\\0,&c\nmid d.\end{cases}
\]
Therefore
\begin{equation}\label{eq:legacy-8.1}
\boxed{E_{m,d}E_{m,e}=E_{m,\gcd(d,e)}.}\end{equation}
M\"obius inversion produces pairwise orthogonal exact-conductor projectors
\begin{equation}\label{eq:legacy-8.2}
\boxed{
Q_{m,c}=\sum_{d\mid c}\mu(c/d)E_{m,d}.
}\end{equation}
They satisfy
\[
Q_{m,c}Q_{m,e}=\delta_{ce}Q_{m,c},
\qquad
\sum_{c\mid m}Q_{m,c}=I.
\]
Hence
\begin{equation}\label{eq:legacy-8.3}
\boxed{
H^4_{\mathrm{prim}}(X_m^4,\Q)
=\bigoplus_{c\mid m}\operatorname{im}Q_{m,c}.
}\end{equation}
Each \(Q_{m,c}\) is a rational linear combination of graph correspondences, so the conductor decomposition is itself algebraic.



\begin{figure}[tbp]\centering\includegraphics[width=0.98\linewidth]{12_conductor_descent.pdf}\caption{Exact conductor descent in the named example. Coordinate-power inflation enlarges the ambient degree while preserving the old conductor.}\end{figure}

\chapter{Algebraic certificate grammar}
Fix a full-conductor orbit in \(\mathcal B_m\). The constructive program tests a hierarchy of increasingly general certificate channels.

\section{Complementary pairs}
For \(1\le r<m\), let
\[
p_r=e_r+e_{m-r}.
\]
A packet is pair-decomposable when
\begin{equation}\label{eq:legacy-9.1}
n=p_{r_1}+p_{r_2}+p_{r_3}.\end{equation}
These are the standard linear-cycle or Tate-pair packets.

\section{Balanced four-packets and saturation}
Let
\[
\mathcal F_m=
\left\{
\gamma\in\mathbf N^{m-1}:
\mathbf1^T\gamma=4,
\quad R_m\gamma=2m\,\mathbf1
\right\}.
\]
A saturated certificate is an exact nonnegative identity
\begin{equation}\label{eq:legacy-9.2}
\boxed{
 n+\sum_{j=1}^{2N-3}p_{r_j}
 =\sum_{i=1}^N\gamma_i,
 \qquad \gamma_i\in\mathcal F_m.
}\end{equation}
The case \(N=2\) is one-step quasi-decomposition. Higher \(N\) allows several surface packets to reconstruct one six-packet after Tate stabilization.

The algebraic implication must use the actual classical correspondence operations. A balanced four-entry packet is algebraic by Lefschetz $(1,1)$ on a Fermat surface. Concatenate the $N$ surface characters using the Shioda--Ran construction, then remove precisely the $2N-3$ complementary pairs by the classical pair-removal theorem. The ambient dimension rises to $4N-2$ before cancellation. This proves algebraicity of a packet admitting such an identity once those cited operations and their character conventions have been checked. It does not authorize cancellation of arbitrary tensor factors. The August 15 status note correctly demanded a complete correspondence-level account of the proposed general grammar. The concrete degree-210 application below spells out the necessary stabilization and uses only the checked pair-removal theorem.

\section{Dual charges}
An antisymmetric charge is a rational function \(w\) on residues satisfying
\[
w(m-r)=-w(r)
\]
and
\[
\sum_r\gamma_r w(r)=0
\qquad\forall \gamma\in\mathcal F_m.
\]
Such a charge vanishes on every complementary pair and on every saturated identity. Writing $w(n)=\sum_r n_r w(r)$, if \(w(n)\ne0\), the entire lower-packet saturation semigroup is excluded. This is a universal obstruction to a grammar, not a nonalgebraicity theorem.

\section{Curve, Delsarte, and incidence channels}
A curve certificate splits the six entries into triples
\[
a=A\sqcup B,
\qquad \sum A=m,
\quad \sum B=2m,
\quad B=uA\pmod m.
\]
The resulting \((1,1)\) class on a product of Fermat curves is algebraic by Lefschetz \((1,1)\), and the join correspondence returns a surface class.

A Delsarte certificate consists of an exponent matrix and a resolved monomial map to a lower-degree hypersurface carrying an already algebraic packet. An incidence certificate uses algebraic lowering maps
\[
A_r:V\to H^2(Y_r)
\]
whose combined positive transfer
\[
T=\sum_rA_r^\dagger A_r
\]
is invertible on the packet. Divisor lowering and Cayley--Hamilton reconstruction then produce the cycle.

\chapter{The degree-33 rejection branch and its collapse}
The degree-33 Fermat fourfold provided a stringent test because one rational orbit escaped the ordinary decomposable and standard channels. A rejection attempt first isolated the packet, excluded several obvious cycle families, and then discovered a cross-family cubic correspondence that algebraizes it.

\section{The exceptional packet}
A representative is
\[
a=(7,10,13,19,22,28),
\]
or after permutation,
\[
\widetilde a=(7,19,28,10,13,22).
\]
Every unit conjugate has residue sum \(99=3\cdot33\), so the orbit is rational of type \((2,2)\). It is not totally pair-decomposable. The previously proposed symmetric surface
\[
W=\{p_1(x^{11})=p_2(x^{11})=p_3(x^{11})=0\}
\]
does not carry the primitive packet: Newton's identity makes the third equation redundant on \(X_{33}^4\), and
\[
[W]=242h^2.
\]
Thus \(W\) is purely ambient.

\section{Klein--Delsarte cubic}
Consider the smooth cubic fourfold
\[
Y:\quad y_0^2y_4+y_1^2y_0+y_2^2y_1+y_3^2y_2+y_4^2y_3+y_5^3=0.
\]
An explicit monomial exponent matrix gives a dominant rational map
\[
\phi:X_{33}^4\dashrightarrow Y
\]
whose projective-torus lattice map has Smith diagonal
\[
(1,33,33,33,33),
\]
so the generic torus degree is \(33^4\).

Take the cubic Jacobian-ring class
\[
P(y)=y_1y_2y_5.
\]
Its exponent vector is $b=(0,1,1,0,0,1)$. Write $C=33A^{-1}$ for the cubic exponent matrix $A$. The class is nonzero in the degree-three Jacobian ring. The residue pullback has the exact exponent identity
\begin{equation}\label{eq:legacy-10.1}
C^Tb+10\mathbf1=(6,18,27,9,12,21)=\widetilde a-\mathbf1.\end{equation}
Therefore the pullback of the corresponding primitive cubic \((2,2)\)-class lies in the one-dimensional Fermat line \(V(\widetilde a)\).

The Laurent exponent match is a useful exact certificate, but by itself it does not prove that a rational class lifts through a compact graph. A complete manuscript argument instead uses the torus-carrier criterion developed in Chapter~\ref{ch:carrier70}. For the reordered character, the exponent matrix $A$ of the cubic satisfies
\[
A^T\widetilde a/33=\beta=(1,2,2,1,1,2),\qquad C=33A^{-1}.
\]
With $r_0$ the first row of $A$, put $u=\beta-3r_0$. Then $\mathbf1^Tu=0$ and $C^Tu=\widetilde a-99e_0$. These identities establish triviality of the selected character on the torus deck group. Every selected character has full support, so its restriction to the coordinate torus is an isomorphism on the packet. Finite étale descent gives a pure rational weight-four Hodge substructure on the cubic torus. The lowest-weight image theorem and semisimplicity give a rational $(2,2)$ lift to the smooth cubic. Classical cubic algebraicity then supplies a surface class. Pull back through a resolved graph, apply the rational orbit projector, and use boundary injectivity to identify the returned class.

\begin{manuscriptclaim}[Degree-33 packet construction]
Subject to the resolved graph and the cited mixed-Hodge inputs in the preceding argument, the rational packet generated by $(7,10,13,19,22,28)$ is algebraic on $X_{33}^4$. The exponent identities are finite exact certificates. The compact geometric implication is a manuscript argument and is not certified by those identities alone.
\end{manuscriptclaim}



The episode is methodologically important. Failure of a finite cycle grammar is not evidence of nonalgebraicity. It may indicate a missing cross-family correspondence.

\begin{figure}[tbp]\centering\includegraphics[width=0.98\linewidth]{15_klein33_route.pdf}\caption{The degree-33 torus/cubic construction requires deck invariance, full-support boundary control, a rational Hodge lift and the resolved algebraic graph.}\end{figure}

\chapter{Infinite towers and full-conductor promotion}
The coordinate-power maps propagate every solved seed through infinitely many degrees.

\begin{theorem}[Infinite packet tower]
Let \(W\) be an algebraic rational character packet on \(X_d^{2p}\). For every \(k\ge1\), its coordinate-power pullback is an algebraic packet on \(X_{kd}^{2p}\). Applied to the preceding degree-33 construction under its geometric hypotheses, this gives a 20-dimensional algebraic packet on every $X_{33k}^4$.
\end{theorem}

The tower packet remains an old conductor-33 sector. A stronger construction appends a complementary pair
\[
(1,m-1)
\]
after scaling a solved seed to modulus \(m=kd\). If
\[
b=(b_0,\ldots,b_{2p+1})
\]
is algebraic at degree \(d\), then
\[
\Theta_{m,b}=(kb_0,\ldots,kb_{2p+1},1,m-1)
\]
has full raw conductor \(m\) and middle Hodge type two dimensions higher. Its rational orbit has dimension \(\varphi(m)\). This produces full-raw-conductor algebraic packets of unbounded rank.

Raw conductor is not the complete novelty invariant: removing Tate pairs or reversing joins can reveal a smaller stable conductor. The stable residual should therefore be defined only after eliminating conductor descent, Tate suspension, joins, Delsarte pullbacks, and incidence channels.

\chapter{Prime-power classification}
The following manuscript classification uses a primitive-character filtration and a six-entry mass bound. At filtration level one the kernel-to-lower-level argument must be replaced by the base-prime Fourier argument stated below.

\section{Valuation strata and orbit mass}
Let \(m=p^r\) with \(p\) odd. Decompose the entries by exact \(p\)-adic valuation. At one stratum, write \(a=p^jb\), with \(b\in U_{p^{r-j}}\), and define the odd count
\[
\delta_j(b)=n_j(b)-n_j(-b).
\]
Since the packet has six entries,
\begin{equation}\label{eq:legacy-12.1}
\|\delta_j\|_1\le 12.\end{equation}
For $s\ge2$, primitive odd characters force \(\delta_j\) to be constant on the kernel cosets of
\[
U_{p^s}\to U_{p^{s-1}}.
\]
A nonzero odd kernel-orbit defect has mass at least \(2p\).

At $s=1$, all odd characters modulo $p$ are primitive. Their sawtooth Fourier coefficients are nonzero, so the odd count vector has every Fourier coefficient zero and is zero. This gives complementary pairing at the base prime level. The higher-level mass argument applies only for $s\ge2$. The needed odd-character nonvanishing is a classical analytic leaf, checked through Sharifi's cyclotomic-field notes \cite{SharifiDirichlet}.

\section{Large odd primes}
If \(p\ge7\), then \(2p>12\), so every odd defect vanishes. Each residue occurs with the same multiplicity as its complement. Therefore every packet is three complementary pairs.

\section{The three-primary branch}
For $p=3$ and $r\ge2$, a higher-level kernel orbit has three elements. Define
\[
O_b=\{b,b+q,b+2q\},
\qquad q=3^{r-1},
\]
and
\[
\Gamma_b=O_b\sqcup\{3(q-b)\}.
\]
Take $1\le b<q$ with $3\nmid b$, ensuring all entries are nonzero. Then \(\Gamma_b\) is a balanced four-packet. A case analysis on the number of oriented kernel orbits gives exact identities expressing every full-conductor nonpair packet as a pair/surface packet or a one-step saturated sum of two balanced four-packets.

\section{The five-primary branch}
For $p=5$ and $r\ge2$, a nonzero odd defect consumes exactly five unit entries. The base case $r=1$ was already handled by the base-prime Fourier argument. The sixth entry is forced by balance. Up to the unit action there is one full-conductor exceptional orbit
\begin{equation}\label{eq:legacy-12.2}
\boxed{
\alpha_{5^r}=
(1,1+q,1+2q,1+3q,1+4q,5q-5),
\qquad q=5^{r-1}.
}\end{equation}
All other packets lie in the pair/surface grammar. The family \eqref{eq:legacy-12.2} is the pure five-primary Hasse--Davenport branch.



\begin{center}
\small
\begin{tabular}{rrrr}\toprule $m$ & Full orbits & Pair/SAT(2) & Residual \\
\midrule
49 & 124 & 124 & 0 \\
81 & 436 & 436 & 0 \\
125 & 827 & 826 pair & 1 \\
\bottomrule\end{tabular}

\end{center}

\chapter{Weighted dyadic reconstruction}
Prime \(2\) cannot be handled by inserting \(p=2\) into the odd-prime orbit-mass proof. The unit group has the separate form
\[
\mathbf Z_2^\times\cong\{\pm1\}\times(1+4\mathbf Z_2).
\]
The first address is a sign packet, only afterward does the principal-unit refinement begin.

Let $s\ge3$, \(m=2^s\), \(h=2^{s-1}\), and write $\rho_s=\rho_{2^s}$. For an odd residue \(b\pmod h\), its two lifts are \(b\) and \(b+h\). Define
\[
L_sn(b)=n(b)+n(b+h),
\qquad
D_sn(b)=n(b)-n(b+h).
\]
For every odd unit \(u\), one has the exact weighted identity
\begin{equation}\label{eq:legacy-13.1}
\boxed{
\begin{aligned}
&n(b)\rho_s(ub)+n(b+h)\rho_s(u(b+h))\\
&\qquad=L_sn(b)\rho_{s-1}(\bar ub)
+2^{s-2}D_sn(b)\varepsilon_s(u,b),
\end{aligned}
}\end{equation}
where \(\varepsilon_s(u,b)\in\{\pm1\}\). Repetition yields a Haar-like expansion whose detail packets carry visible powers of two. The Hodge balance equation becomes a binary carry relation between weighted details.



Exact enumeration gives:
\begin{center}
\small
\begin{tabular}{rrrrrr}\toprule $m$ & Full orbits & Pair & SAT(2) & SAT(3) & Unresolved \\
\midrule
8 & 13 & 9 & 4 & 0 & 0 \\
16 & 45 & 27 & 18 & 0 & 0 \\
32 & 134 & 91 & 42 & 1 & 0 \\
64 & 419 & 331 & 87 & 1 & 0 \\
128 & 1435 & 1259 & 175 & 1 & 0 \\
\bottomrule\end{tabular}

\end{center}
The table suggests one residual orbit beginning at \(m=32\). The remaining derivation proposes an all-degree explanation. Its uniform status was downgraded by the later August 15 note, so this edition preserves the argument as a manuscript claim requiring a complete audit.

\begin{figure}[tbp]\centering\includegraphics[width=0.98\linewidth]{13_dyadic_sat3.pdf}\caption{The three finite dyadic representatives have exact SAT(3) identities. These identities do not prove all-degree grammar coverage.}\end{figure}

\section{Complement defects and dyadic triangles}
Put
\[
q=\frac m4,
\qquad h=2q,
\qquad m=4q.
\]
For a packet count vector \(n\), define its oriented complement defect
\[
\partial n=\sum_{r=1}^{h-1}d_n(r)[r],
\qquad
d_n(r)=n_r-n_{m-r}.
\]
The self-complementary residue \(h\) contributes zero. Since
\(\rho_m(m-r)=-\rho_m(r)\), the balanced-packet equation is equivalent to
\[
\sum_{r=1}^{h-1}d_n(r)\rho_m(ur)=0
\qquad\text{for every odd unit }u.
\]
Moreover,
\[
\|\partial n\|_1=\sum_{r=1}^{h-1}|d_n(r)|\le 6.
\]
Thus the dyadic problem is a short-vector problem in an integral defect lattice.

For \(1\le b<q\), define the triangle defect
\[
\kappa_b=[2b]-[b]+[h-b].
\]
It is realized by the four-packet
\[
\Gamma_b=(2b,h-b,h,m-b).
\]

\begin{manuscriptclaim}[Dyadic triangle relation]
For every odd unit \(u\) and every \(1\le b<q\),
\[
\rho_m(2ub)-\rho_m(ub)+\rho_m(u(h-b))=0.
\]
Consequently \(\kappa_b\) is balanced and \(\Gamma_b\in\mathcal F_m\).
\end{manuscriptclaim}
\begin{proof}
Let \(x=\langle ub\rangle_m\). If \(0<x<h\), then
\[
\rho_m(2x)=2x-h,
\qquad
\rho_m(x)=x-h,
\qquad
\rho_m(h-x)=-x.
\]
Their indicated combination is zero. If \(h<x<m\), write \(x=m-y\) with \(0<y<h\), then
\[
\rho_m(2x)=h-2y,
\qquad
\rho_m(x)=h-y,
\qquad
\rho_m(h-x)=y,
\]
and the same cancellation occurs. Since \(u\) is odd, \(uh\equiv h\pmod m\), so \(u(h-b)\equiv h-ub\pmod m\). This proves the identity.
\end{proof}

\begin{figure}[tbp]\centering\includegraphics[width=0.98\linewidth]{14_dyadic_tree32.pdf}\caption{The finite degree-32 defect tree. The proposed uniform derivation is retained separately with its publication-status boundary.}\end{figure}

\section{The integral triangle basis}
Let
\[
\mathcal K_m=
\left\{
 d\in\mathbf Z^{h-1}:
 \sum_{r=1}^{h-1}d(r)\rho_m(ur)=0
 \text{ for every odd }u
\right\}.
\]
The triangle relations do not merely generate examples, they form the whole balanced defect lattice.

\begin{manuscriptclaim}[Integral dyadic triangle basis]
For \(m=2^s\), \(s\ge3\), the vectors
\[
\{\kappa_b:1\le b<q\}
\]
form a \(\mathbf Z\)-basis of \(\mathcal K_m\).
\end{manuscriptclaim}
\begin{proof}
Set \(r_0=q/2\). On the vertex set
\[
V_m=\{1,2,\ldots,q-1\},
\]
define, for \(b\ne r_0\),
\[
p(b)=
\begin{cases}
2b,&b<r_0,\\
2(q-b),&b>r_0,
\end{cases}
\qquad
\epsilon_b=
\begin{cases}
+1,&b<r_0,\\
-1,&b>r_0.
\end{cases}
\]
The map \(p\) raises the two-adic valuation by one, so repeated application reaches \(r_0\). It therefore defines a rooted tree.

Suppose \(d=\sum_{b=1}^{q-1}z_b\kappa_b\). Its even coordinates are
\[
d(q)=z_{r_0},
\qquad
d(2b)=z_b-\epsilon_bz_{p(b)}
\quad(b\ne r_0).
\]
Starting at the root and moving toward the leaves, these equations recover every \(z_b\) integrally. Hence the triangle vectors are linearly independent and their even-coordinate map is unimodular.

Now take \(d\in\mathcal K_m\), recover the unique integers \(z_b\) from its even coordinates, and subtract \(\sum_bz_b\kappa_b\). The residual balanced vector is supported only on the odd residues \(1\le r<h\). The corresponding square sawtooth matrix is nonsingular. Indeed, extend an odd-support vector to an odd function on \(U_m\). Convolution with \(\rho_m\) is diagonalized by the odd Dirichlet characters, with eigenvalues
\[
\lambda_\chi=
\sum_{a\in U_m}\rho_m(a)\overline{\chi(a)}
=
\sum_{a\in U_m}a\overline{\chi(a)}.
\]
If \(\chi\) is induced from a primitive odd character \(\chi^*\) of conductor \(2^t\), then
\[
\lambda_{\chi,2^s}=2^{s-t}\lambda_{\chi^*,2^t}.
\]
The primitive value is, up to a nonzero scalar, the generalized Bernoulli value \(B_{1,\overline{\chi^*}}\), equivalently \(-L(0,\overline{\chi^*})\), and is nonzero for every odd primitive character \cite{WashingtonCyclotomic,SharifiDirichlet}. Thus every odd Fourier eigenvalue is nonzero. The odd residual must vanish, so the \(\kappa_b\) span \(\mathcal K_m\) over \(\mathbf Z\).
\end{proof}

\section{The trivalent defect tree}
The integral basis turns defect mass into exact tree variation. Define a sign gauge by
\[
\tau_{r_0}=1,
\qquad
\tau_b=\epsilon_b\tau_{p(b)},
\qquad
y_b=\tau_bz_b.
\]
Attach one zero-valued boundary leaf to the root and two zero-valued boundary leaves to every odd vertex. Every vertex of the resulting augmented tree \(\widetilde T_m\) has degree three. The boundary edges correspond bijectively to defect coordinates: the root edge is labelled \(q\), the edge from \(b\) to \(p(b)\) is labelled \(2b\), and the two leaves at an odd \(b\) are labelled \(b\) and \(h-b\).

\begin{manuscriptclaim}[Exact trivalent-tree norm]
If \(d=\sum_bz_b\kappa_b\in\mathcal K_m\), then
\begin{equation}\label{eq:legacy-13.2}
\boxed{
\|d\|_1
=|y_{r_0}|
+\sum_{b\ne r_0}|y_b-y_{p(b)}|
+2\sum_{\substack{1\le b<q\\ b\text{ odd}}}|y_b|.
}\end{equation}
The right side is the total variation of the integer function \(y\) on \(\widetilde T_m\) with zero boundary values.
\end{manuscriptclaim}
\begin{proof}
The root coordinate is \(d(q)=z_{r_0}\). For \(b\ne r_0\), multiplying
\(d(2b)=z_b-\epsilon_bz_{p(b)}\) by \(\tau_b\) gives
\[
\tau_bd(2b)=y_b-y_{p(b)}.
\]
For odd \(b\), the two odd coordinates are
\[
d(b)=-z_b,
\qquad
d(h-b)=z_b.
\]
Taking absolute values and summing all defect coordinates gives \eqref{eq:legacy-13.2}.
\end{proof}

For an integer-valued \(y\), discrete coarea gives
\begin{equation}\label{eq:legacy-13.3}
\operatorname{TV}(y)
=
\sum_{k\ge1}|\partial\{y\ge k\}|
+
\sum_{k\ge1}|\partial\{y\le-k\}|.\end{equation}
If \(S\) is a finite forest of internal vertices of a trivalent tree and \(c(S)\) is its number of connected components, then
\begin{equation}\label{eq:legacy-13.4}
|\partial S|=|S|+2c(S),\end{equation}
because \(3|S|=2(|S|-c(S))+|\partial S|\).

\begin{manuscriptclaim}[All defect shapes of mass at most six]
Let \(d\in\mathcal K_m\) satisfy \(\|d\|_1\le6\), and let \(y\) be its gauged coefficient function. Exactly one of the following occurs:
\begin{enumerate}
\item \(y=0\)
\item \(y=\pm2\) at one vertex and is zero elsewhere
\item \(y\) has two singleton components, each with value \(\pm1\)
\item \(y=\pm1\) on one connected subtree containing at most four vertices and is zero elsewhere.
\end{enumerate}
For a balanced four-packet, only the zero shape, a single vertex, or one tree edge can occur.
\end{manuscriptclaim}
\begin{proof}
Apply \eqref{eq:legacy-13.3}--\eqref{eq:legacy-13.4}. A nonempty level-set component with \(k\) vertices costs \(k+2\). If \(|y|\ge2\), two nonempty nested level sets already cost at least six, equality forces both to be the same singleton. If positive and negative values both occur, their first level sets cost at least three each, equality forces two singleton components. Otherwise \(y\) has one sign, and a connected support of size \(k\) costs \(k+2\), so \(k\le4\). Two same-sign components can have total cost at most six only when both are singletons. Replacing six by four leaves only a vertex or an edge.
\end{proof}

\section{SAT(2) closure and the four-star obstruction}
For a defect \(d\), let \(\operatorname{red}(d)\) be the unmatched residue multiset: a positive coefficient at \(r\) contributes copies of \(r\), while a negative coefficient contributes copies of \(m-r\). Its length is \(\|d\|_1\). Every packet with defect \(d\) is \(\operatorname{red}(d)\) together with complementary pairs and copies of \(h\).

A one-vertex tree atom has defect \(\pm\kappa_b\) after the gauge sign and is represented by \(\Gamma_b\) or its conjugate, its reduced length is three and its fourth entry is \(h\). A one-edge atom has defect mass four, so its reduced multiset itself is a balanced four-packet. A zero atom is any union of two complementary pairs.

\begin{manuscriptclaim}[Exact SAT(2) criterion]
Let \(n\) be a balanced six-packet at modulus \(m=2^s\). Then \(n\) is pair-decomposable or has an SAT(2) certificate unless its gauged defect is \(\pm1\) on a four-vertex star \(K_{1,3}\) in the internal tree \(T_m\). A four-star has no SAT(2) certificate.
\end{manuscriptclaim}
\begin{proof}
Every nonstar shape in the preceding proposition splits into at most two four-packet atoms:
\begin{itemize}
\item a vertex or edge is paired with a zero atom
\item a three-vertex path is an edge plus its remaining terminal vertex
\item a four-vertex path is the union of two disjoint edges
\item a doubled singleton is two copies of one vertex atom
\item two singleton components are the corresponding two vertex atoms.
\end{itemize}
These defect splittings lift to exact nonnegative multiset identities. For a vertex, the neutral remainder of the six-packet is \(h\) plus one complementary pair, for an edge it is one complementary pair. A zero four-packet absorbs that pair and the stabilizing pair. Cutting a three- or four-vertex path produces exactly one oppositely oriented boundary edge, hence exactly one complementary pair in the union of the two atom packets, this is the SAT(2) stabilizer. In the three-vertex case, the vertex atom also supplies the unique neutral \(h\). For two vertex atoms, their two neutral entries are \((h,h)=p_h\). The zero-defect case is already a union of three complementary pairs. Thus in every nonstar case there are \(\gamma_1,\gamma_2\in\mathcal F_m\) and a complementary pair \(p_r\) such that
\[
n+p_r=\gamma_1+\gamma_2.
\]

Conversely, suppose a four-star defect had an SAT(2) certificate. Each balanced four-packet has tree support equal to zero, one vertex, or one edge. The sum of two such supports must produce four nonzero star vertices. Therefore both supports must be disjoint edges covering all four vertices. This would be a perfect matching of \(K_{1,3}\), which does not exist. Hence the star is the unique tree obstruction to SAT(2).
\end{proof}

\section{Conductor of a star and its unique unit orbit}
The tree also determines exactly when the star obstruction has full conductor.

\begin{manuscriptclaim}[Star conductor]
Let a four-star have center \(c\), and write \(j=v_2(c)\). The gcd of \(m\) with its six packet entries is
\[
2^{j-1}.
\]
Hence a star has full conductor exactly when \(v_2(c)=1\).
\end{manuscriptclaim}
\begin{proof}
Use the boundary-edge labels of \(\widetilde T_m\). The two children of \(c\) have valuation \(j-1\). The four boundary edges leaving those children have labels of exact valuation \(j-1\): for \(j=1\) they are the odd leaf labels, and for \(j>1\) they are labels \(2b\) coming from grandchildren of valuation \(j-2\). The two boundary edges leaving the parent side have larger valuation. Thus every nonzero defect coordinate is divisible by \(2^{j-1}\), and at least four have exact valuation \(j-1\). Replacing a defect coordinate \(r\) by the packet entry \(m-r\) does not change its two-adic valuation. The stated gcd follows.
\end{proof}

Assume now \(s\ge5\) and let a full-conductor star have center \(c=2a\) with \(a\) odd. Up to global conjugation, expanding its four triangle atoms gives the unordered packet
\begin{equation}\label{eq:legacy-13.5}
P_a=
(8a,q-a,2q-a,2q-4a,3q-a,4q-a).\end{equation}
The cases \(c<q/2\) and \(c>q/2\) give the same unordered tuple after the internal-edge cancellations. Choose an odd \(u\) satisfying
\[
ua\equiv1\pmod{2q}.
\]
Since \(uq\equiv q\) or \(3q\pmod{4q}\), multiplication by \(u\) sends \eqref{eq:legacy-13.5} to
\[
P_1=(8,q-1,2q-4,2q-1,3q-1,4q-1).
\]
Its conjugate is
\begin{equation}\label{eq:legacy-13.6}
\boxed{
A_m=
\left(
1,\frac m4+1,\frac m2+1,\frac m2+4,
\frac{3m}4+1,m-8
\right).
}\end{equation}
Thus every full-conductor star belongs to one \(U_m\)-orbit.

The packet \(A_m\) obeys the exact uniform SAT(3) identity
\begin{equation}\label{eq:legacy-13.7}
\begin{aligned}
&A_m+(2,m-2)+(4,m-4)+(m/2,m/2)\\
&=(1,m/2,m/2+1,m-2)\\
&\quad +(2,m/4+1,3m/4+1,m-4)\\
&\quad +(4,m/2,m/2+4,m-8).
\end{aligned}\end{equation}
The first and third packets on the right have defects \(-\kappa_1\) and \(-\kappa_4\). The middle packet has defect \(-\kappa_2+\kappa_{q-1}\), an edge atom. Hence all three are balanced four-packets, and \eqref{eq:legacy-13.7} is an identity of multisets.

\begin{manuscriptclaim}[All-degree dyadic uniqueness]
Let \(m=2^s\), \(s\ge3\). Every full-conductor balanced six-packet is pair-decomposable or SAT(2), except that for \(s\ge5\) there is exactly one residual unit orbit, represented by \(A_m\). This orbit is SAT(3) by \eqref{eq:legacy-13.7}. For \(m=8\) and \(m=16\), no four-star exists and there is no residual orbit.
\end{manuscriptclaim}
\begin{proof}
The trivalent-tree norm and coarea classification list every defect of mass at most six. The SAT(2) criterion closes every listed shape except a four-star. The star-conductor lemma forces a full-conductor star center to have valuation one, and the calculation \eqref{eq:legacy-13.5}--\eqref{eq:legacy-13.6} places all such stars in the single orbit of \(A_m\). The trees for \(m=8\) and \(m=16\) have no internal vertex with three internal neighbors. Finally, \eqref{eq:legacy-13.7} closes the unique residual orbit at saturation depth three.
\end{proof}

\begin{manuscriptclaim}[Complete dyadic grammar]
After exact conductor descent, every balanced packet at every dyadic modulus is generated by complementary pairs, SAT(2) surface packets, and coordinate-power pullbacks of the single SAT(3) family \(A_{2^t}\). No second residual dyadic family occurs at any depth.
\end{manuscriptclaim}

\begin{resultbox}[Uniform dyadic status]
The August 6 manuscript called uniform uniqueness closed. The August 15 status note restored it to a uniform write-up obligation. The proposed proof uses an integral triangle basis, an exact total-variation model, discrete coarea, and star conductor. This edition does not reverse the later downgrade on the strength of finite agreement.
\end{resultbox}

The archived companion dyadic verification report records complete orbit counts through $m=128$ and structural identities through $2^{14}$. The October replay freshly reconstructs the complete census through 128 and the triangle/tree identities through $2^9$. The report file is preserved evidence, rather than an executable verifier or a current replay of the larger historical range. These computations support the finite identities. They do not supply a current independent proof of the all-degree claim.

\chapter{Finite global audit through degree 100}
The finite audit exactly enumerates balanced packets, reduces them by unit orbits, checks pair and SAT(2) certificates, and searches higher saturated identities with \(N\le8\). Every returned identity is rechecked by exact integer arithmetic.

The archived named verification groups report the following outcomes:
\begin{center}
\small
\begin{tabular}{lrr}\toprule Suite & Passed & Total \\
\midrule
projector resolved hodge calculus & 18 & 18 \\
fermat33 klein closure & 14 & 14 \\
infinite klein tower & 4 & 4 \\
conductor polyphase & 4 & 4 \\
odd prime power audit & 13 & 13 \\
\midrule Total & 53 & 53 \\
\bottomrule\end{tabular}

\end{center}
The archived dyadic report records \(89{,}478{,}480\) coefficient checks through \(2^{14}\). The October replay independently checked 87,376 coefficients for exponents $3\le s\le9$. The larger historical count was not rerun in full. The archived packet audit reports completeness through degree \(100\) for its implemented grammar, which is not a completeness theorem for all correspondence channels.

After removing the degree-33 Klein channel and the degree-39 curve channel, every remaining post-saturation degree in the audit is divisible by five. This is a powerful localization of the residual grammar, but it is not by itself a universal proof.



\section{Exact residual representatives}
\small
\begin{longtable}{rp{0.68\textwidth}}\toprule Degree & Post-saturation representatives \\
\midrule\endfirsthead
\toprule Degree & Post-saturation representatives \\
\midrule\endhead
25 & \texttt{(1,6,11,16,20,21)} \\
33 & \texttt{(1,4,16,22,25,31)} \\
35 & \texttt{(1,8,15,22,29,30)} \\
39 & \texttt{(1,7,16,22,34,37)} \\
40 & \texttt{(1,9,17,25,33,35)} \\
45 & \texttt{(1,10,19,28,37,40); (1,19,20,28,30,37)} \\
50 & \texttt{(1,7,27,30,41,44); (1,11,21,31,41,45); (1,13,22,32,40,42)} \\
55 & \texttt{(1,12,23,34,45,50)} \\
65 & \texttt{(1,14,27,40,53,60)} \\
70 & \texttt{(1,6,43,45,57,58); (1,15,29,43,57,65); (1,20,24,42,61,62); (1,20,29,43,57,60); (1,29,35,43,45,57)} \\
75 & \texttt{(1,16,31,46,61,70)} \\
80 & \texttt{(1,17,33,49,65,75)} \\
85 & \texttt{(1,18,35,52,69,80)} \\
90 & \texttt{(1,8,26,62,85,88); (1,19,37,55,73,85); (1,19,37,60,73,80); (1,19,37,65,73,75); (1,23,38,55,74,79); (1,38,46,55,56,74); (2,38,45,55,56,74)} \\
95 & \texttt{(1,20,39,58,77,90)} \\
100 & \texttt{(1,21,41,61,81,95)} \\
\bottomrule\end{longtable}

\normalsize

The residuals are not random. Their mod-5 phase profiles form a small local grammar.

\compactpart{Part VIII. CRT--Kummer Polyphase and the Imaginary CM Half-Twist}

\chapter{Local mod-5 phase grammar}
Assume \(m=5q\) and first suppose \((5,q)=1\). For \(b\pmod q\) and a phase \(r\in\{0,1,2,3,4\}\), let \(\lambda_r(b)\) be the CRT lift satisfying
\[
\lambda_r(b)\equiv r\pmod5,
\qquad
\lambda_r(b)\equiv b\pmod q.
\]
The finite audit reveals four canonical multiplicity profiles, up to multiplication by a unit modulo five:
\[
(1,0,0,0,5),
\quad
(2,1,1,1,1),
\quad
(1,0,2,1,2),
\quad
(1,0,1,3,1).
\]
The first is the complete-five-phase K5 packet, the second is a four-phase K4 lift, and the final two are one K23 pair--triple mechanism.



\section{K4 four-phase reduction}
Choose nonzero residues \(b,c,d\pmod q\) and put
\[
e\equiv-5^{-1}b\pmod q.
\]
Define
\[
A_{b,c,d}=
(\lambda_1(b),\lambda_2(b),\lambda_3(b),\lambda_4(b),5c,5d)
\]
and
\[
\Gamma_{b,c,d}=(b,c,d,e).
\]

\begin{theorem}[Four-phase CRT reduction]
The six-packet \(A_{b,c,d}\) is balanced modulo \(5q\) if and only if the four-packet \(\Gamma_{b,c,d}\) is balanced modulo \(q\).
\end{theorem}
\begin{proof}
For a unit \(u\), let \(s=\langle ub\rangle_q\) and \(z=\langle u5^{-1}b\rangle_q\). The complete five-element fiber over \(s\) has sum \(5s+10q\), its phase-zero lift is \(5z\). The four nonzero phases therefore sum to \(5s+10q-5z\). Adding the last two entries and imposing the middle value \(15q\) gives
\[
s+\langle uc\rangle_q+\langle ud\rangle_q-z=q.
\]
Since \(\langle ue\rangle_q=q-z\), this is exactly the balance equation for \((b,c,d,e)\).
\end{proof}

All 16 audited K4 residuals reduce exactly to balanced four-packets.

\chapter{K23 pair--triple factorization}
Let a K23 packet have one phase-zero entry and five nonzero phases:
\[
A=P\sqcup T\sqcup\{5e\},
\qquad |P|=2,
\quad |T|=3,
\]
with
\[
\sum P\equiv0\pmod5,
\qquad
\sum T\equiv0\pmod5.
\]
Write
\[
\sum P=5C,
\qquad
\sum T=5D,
\]
and reduce \(C,D\) modulo \(q\) to representatives \(c,d\). Require $c,d,e$ to be nonzero modulo $q$ so that all three resulting character blocks have full support. Degenerate cases must leave this construction and be handled through pair removal or another channel. For example $(1,14,2,9,14,5)$ at modulus 15 is balanced and admits the indicated phase split, but $P=(1,14)$ gives $c=0$. The phase condition alone is insufficient. Define
\[
C_P=P\sqcup\{m-5c\},
\]
\[
S_T=T\sqcup\{m-5d\},
\]
and
\[
L=(5c,5d,5e).
\]

For any nonzero-entry block $B$ at modulus $m$, put $h_B(u)=m^{-1}\sum_{r\in B}\langle ur\rangle_m$, with multiplicities retained.

\begin{theorem}[Exact K23 packet identity]
One has the exact multiset equality
\begin{equation}\label{eq:legacy-15.1}
\boxed{
A+(5c,m-5c)+(5d,m-5d)=C_P+S_T+L.
}\end{equation}
For every unit \(u\pmod m\),
\[
h_{C_P}(u)+h_{S_T}(u)+h_L(u)=5,
\]
and \(c+d+e\equiv0\pmod q\).
\end{theorem}

The five nonzero phases have only two irreducible possibilities: all five phases are equal, or they contain a complementary pair and a zero-sum triple. Indeed, if no complementary phase pair occurs, the phase support lies in one of four two-element subsets, the zero-sum condition forces all five entries to use one phase. This proves the local K5/K23 dichotomy.



After fixing a common cyclotomic coefficient field and a diagonal character projector, write the selected factor schematically as
\[
N=e_{\mathrm{diag}}\bigl(M(C_P)\otimes M(S_T)\bigr).
\]
An ordinary tensor product of rational packets contains off-diagonal embedding pairs and is generally not rank one. Rank one here refers only to a selected eigenspace after scalar extension. The coefficient field and projector are part of the unresolved correspondence specification.
The Hodge-level identity implies that \(N(1)\) has weight one and has Hodge type opposite to the lower curve packet \(M(L)\) at every embedding.

\begin{figure}[tbp]\centering\includegraphics[width=0.98\linewidth]{16_k23_factorization50.pdf}\caption{A valid finite K23 factorization at modulus 50, with all required character entries nonzero.}\end{figure}

\chapter{Imaginary CM selector}
Write the lower curve packet as
\[
\ell=(c,d,e)\pmod q,
\qquad c+d+e\equiv0\pmod q.
\]
Let \(K=\Q(\zeta_q)\) and define
\begin{equation}\label{eq:legacy-16.1}
\boxed{
\Delta_\ell=(1-\zeta_q^c)(1-\zeta_q^d)(1-\zeta_q^e).
}\end{equation}

\begin{theorem}[Pure imaginary selector]
\[
\overline{\Delta_\ell}=-\Delta_\ell.
\]
\end{theorem}
\begin{proof}
Since
\[
1-\zeta_q^{-a}=-\zeta_q^{-a}(1-\zeta_q^a),
\]
complex conjugation multiplies the product by
\[
(-1)^3\zeta_q^{-(c+d+e)}=-1.
\]
\end{proof}

For an embedding \(\sigma_u:\zeta_q\mapsto\zeta_q^u\), let
\[
h_\ell(u)=\frac{\langle uc\rangle_q+\langle ud\rangle_q+\langle ue\rangle_q}{q}\in\{1,2\}.
\]
Using
\[
1-e^{2\pi ir/q}=-2i\,e^{\pi ir/q}\sin\frac{\pi r}{q},
\]
one obtains
\begin{equation}\label{eq:legacy-16.2}
\sigma_u(\Delta_\ell)
=8i(-1)^{h_\ell(u)}
\prod_{j=1}^3
\sin\frac{\pi\langle u\ell_j\rangle_q}{q}.\end{equation}
Here $(\ell_1,\ell_2,\ell_3)=(c,d,e)$ is an indexed triple, so repeated residues contribute repeated factors. Every sine factor is positive, so the sign of the imaginary part is exactly the Hodge sign. Consequently the CM type of \(N(1)\) is the conjugate of the CM type of the lower curve packet:
\begin{equation}\label{eq:legacy-16.3}
\boxed{
\Sigma_N=\overline{\Sigma_\ell}.
}\end{equation}
This is the rigorous form of the statement that the half-twist occurs on the imaginary plane.



The exact audit generates 364 K23 certificates from seven rational packet orbits. Every certificate obeys the opposite curve CM sign, factors through a lower conductor, and is a Galois translate of an upper-half-plane CM type.

\begin{center}
\small
\begin{tabular}{lr}\toprule Item & Value \\
\midrule
K23 rational orbits & 7 \\
Exact K23 certificates & 364 \\
Opposite curve CM sign & True \\
Lower-conductor factorization & True \\
Upper-half-plane CM translate & True \\
\bottomrule\end{tabular}

\end{center}

In the audited cases, the minimal CM conductors are \(3,5,7,\) or \(9\), even when the original fourfold modulus is \(50,70,\) or \(90\). Thus the weight-one carrier is not arbitrary: it is a character factor of the Jacobian of a lower-degree Fermat curve.

\begin{figure}[tbp]\centering\includegraphics[width=0.98\linewidth]{17_cm_imaginary_sign.pdf}\caption{The exact imaginary selector sign in a retained example. Matching signs does not construct a rational Chow correspondence.}\end{figure}

\chapter{Explicit carrier and remaining Chow map}
For an audited CM type, the lower Fermat curve provides a candidate rational character-orbit factor of its Jacobian. Write $e_{\mathrm{orb}}J(F_d)$ for that factor, up to isogeny. The idempotent selects the complete rational Galois orbit. A projector onto just one holomorphic CM half-plane would not be rational and is not an algebraic rational packet projector.

The desired correspondence is schematically
\[
\Gamma\in\CH^{\ast}(F_d\times C_P\times S_T)_{\Q},\qquad
\Gamma_*:e_{\mathrm{orb}}H^1(F_d)\longrightarrow
e_{\mathrm{diag}}\bigl(M(C_P)\otimes M(S_T)\bigr)(1).
\]
This formula is a target specification, not a constructed map. Its coefficient field, diagonal projector, rational multiplicities, and dimensions must be fixed before an isomorphism can be asserted. The smaller CM conductor can have a smaller rational degree than the original packet. Matching sign functions alone does not resolve that difference.

After the typing has been supplied, a nonzero equivariant map between corresponding rank-one eigenspaces is an isomorphism over the splitting field. Proving that the map is induced by an algebraic correspondence remains essential. One needs a Kummer or half-twist correspondence with the correct covariance and a nonzero period or intersection coefficient. Tensoring a verified map with the conjugate lower curve packet can then place a class in divisor degree on a product of curves or abelian varieties. Lefschetz $(1,1)$ and actual Fermat join maps would transport it back.

\begin{openproblem}[Typed CM packet correspondence]
Specify the fields, projectors and rational multiplicities, construct the algebraic map, and prove its nonzero character coefficient. The 364 audited sign certificates identify candidate factors. They do not solve this problem. The finite-field Chow-motive constructions of Otsubo and Yamazaki provide related context, and are not a characteristic-zero correspondence proof for these packets \cite{OtsuboYamazaki}.
\end{openproblem}

\begin{figure}[tbp]\centering\includegraphics[width=0.98\linewidth]{18_missing_chow_bridge.pdf}\caption{The typed algebraic CM bridge is an open construction. Coefficient field, projectors, multiplicities and nonzero period are part of the target.}\end{figure}

\chapter{Correct literature and novelty map}
Kang is cited in the checked later literature for the all-degree Fermat-fourfold Hodge result, while its original exact theorem and proof were not extracted here \cite{KangFermat}. Separately, da Silva studies which packets are generated by Shioda-style induced-cycle criteria and records the degree-$33$ failure of that criterion \cite{DaSilvaFermat}. The October source contracts authenticate the Kang attribution through the checked later literature while recording that its original proof was not extracted. No historical priority claim is made for this constructive program.

The constructive contribution of the packet program is therefore not best described as the first proof of the rational Hodge conjecture for all Fermat fourfolds. Its strongest defensible targets are:
\begin{itemize}
\item explicit conductor projectors as algebraic correspondences
\item an algorithmic packet certificate grammar
\item a worked cross-family Klein--Delsarte construction at degree 33
\item infinite conductor towers and full-raw-conductor promotions
\item odd prime-power packet classifications
\item a weighted dyadic recursion and a proposed all-degree classification with one SAT(3) orbit
\item exact CRT--Kummer reductions for K4 and K23 packets
\item an imaginary CM selector and explicit lower Fermat-curve carriers
\item an explicit packet-resolved Chow correspondence algorithm
\end{itemize}

\begin{resultbox}[Constructive status]
The abstract existence question for Fermat fourfolds is treated as prior published work. The packet program supplies a detailed constructive architecture and several explicit manuscript correspondence constructions, without a historical novelty claim. Universal constructive coverage also requires a complete mixed-conductor grammar. The typed CM correspondence is one focused gap inside that larger problem.
\end{resultbox}

\compactpart{Part IX. Verification, Audit, and Final Constructive Theorem}

\chapter{Verification ledger}
The complete edition separates symbolic theorems, finite exact checks, geometric correspondences, and literature claims.

\section{Finite checks}
The original Hodge-calculus release passed 18 finite checks. The archived Fermat extension reports the following historical checks, whose fresh replay ranges are distinguished in the current evidence chapter:
\begin{itemize}
\item 14 exact degree-33 Klein checks
\item 4 infinite tower checks
\item 4 exact conductor-polyphase checks
\item 13 odd prime-power checks
\item 89,478,480 dyadic weighted coefficient identities
\item exact packet classifications at dyadic degrees through 128
\item an independent trivalent-tree/SAT(2) audit and structural dyadic checks through $2^{14}$
\item exact balanced-packet audit through degree 100
\item 364 K23 CM-type certificates
\end{itemize}

\section{What finite checks establish}
They establish exact arithmetic identities, matrix products, orbit counts, conductor decompositions, Hodge balance, and nonnegative multiset equalities for every tested instance. They also detect false universal conjectures quickly. The proposed dyadic uniqueness statement is not inferred from these checks. Its manuscript derivation is preserved with the later uniform-write-up downgrade.

\section{What finite checks do not establish}
They do not replace:
\begin{itemize}
\item a universal mixed-prime phase-classification proof beyond the audited range
\item a Chow-level Kummer--Schur or half-twist correspondence
\item a historical priority audit
\end{itemize}

\chapter{Final theorem stack}
The constructive Fermat program can be stated as a dependency stack.

\begin{theorem}[Conditional explicit packet theorem]
Fix a packet that is admitted by the implemented certificate grammar. Assume the following correspondence statements for its chosen channel:
\begin{enumerate}
\item saturated Fermat joins are associative on Chow motives and complementary Tate factors cancel
\item every Delsarte torus map used by the certificate grammar extends through a resolved graph with full-support boundary annihilation
\item every K5, K4, and K23 packet used by the chosen certificate has its correctly typed Chow correspondence with a proved nonzero character coefficient
\end{enumerate}
Then that packet has an explicit packet-resolved algebraic cycle representative. A conclusion for every Fermat fourfold packet additionally requires universal mixed-conductor coverage by this grammar, which has not been proved here.
\end{theorem}
\begin{proof}
Apply the exact conductor projectors. Descend each nonzero component to its full-conductor seed. Use the supplied grammar certificate for the seed, whether pair, saturation, curve, Delsarte, dyadic, or a Kummer phase channel. Use the corresponding Chow map to produce a nonzero cycle projection. Pull old packets back to the original degree and sum the orthogonal conductor components. Finally use the rational orbit projector or rational group-module reconstruction. A bare sum of all complex conjugate eigenvectors may vanish and is not a substitute for that rational construction.
\end{proof}

\begin{openproblem}[Unconditional constructive completion]
Write the missing correspondence maps explicitly and verify their nonzero packet matrix coefficients. A concentrated target is the correctly typed CM correspondence specified above. Universal mixed-conductor coverage is a separate obligation.
\end{openproblem}

\chapter{Research checklist}
\begin{longtable}{p{0.42\textwidth}p{0.17\textwidth}p{0.32\textwidth}}
\toprule
Task & Status & Required proof object \\
\midrule\endfirsthead
\toprule
Task & Status & Required proof object \\
\midrule\endhead
Universal balanced packet equation & Complete & Finite residue matrix \\
Exact conductor descent & Complete & Coordinate-power graph \\
M\"obius conductor projectors & Complete & Rational graph combination \\
Degree-33 Klein packet & Manuscript-level complete & Resolved graph and residue pullback \\
Infinite old-packet towers & Complete & Pullback/pushforward identity \\
Odd prime-power classification & Symbolically complete & Primitive-character filtration \\
Dyadic finite identities & Replayed in a bounded range & Weighted coefficient identities and finite tree data \\
K4 CRT reduction & Complete & Lower balanced four-packet \\
K23 local phase reduction & Audited instances & Pair--triple identity with nonzero-entry conditions \\
K23 CM sign descent & Complete & Imaginary selector $\Delta_\ell$ \\
Explicit lower CM carrier & Identified & Fermat-curve character idempotent \\
CM carrier Chow map & Open & Nonzero algebraic correspondence \\
All-degree constructive algorithm & Open & Universal coverage plus all required correspondence maps \\
Historical priority map & Open audit & Compare existence and constructive literature \\
\bottomrule
\end{longtable}

\chapter{Complete-edition conclusion}
The projector-resolved approach transforms the Hodge problem into a hierarchy of exact finite packets and explicit geometric transfers. For arbitrary varieties, the master obstruction remains the Hodge packet not occupied by algebraic-cycle channels. For Fermat fourfolds, diagonal symmetry makes every stage concrete: balanced residue equations, conductor projectors, prime-local refinements, nonnegative packet identities, and rank-one character motives.

The degree-33 rejection attempt illustrates the central lesson. A packet may evade every obvious lower-packet construction and still be algebraic through a hidden cross-family correspondence. A completed Klein--Delsarte correspondence supplies a seed for infinite conductor towers. The odd prime-power manuscript classification and the proposed dyadic tree derivation explain why local primes require distinct packet grammars. The audited mixed-prime residuals display a small Kummer phase system, without a proof that these patterns exhaust arbitrary conductors. Finally, the imaginary CM selector identifies the precise upper/lower embedding split and reduces the unknown weight-one carrier to a small Fermat-curve Jacobian factor.

The constructive endpoint is now sharply defined:
\begin{resultbox}[Final constructive endpoint]
Construct a correctly typed nonzero equivariant Chow correspondence for each admitted CM packet, and separately prove universal coverage of the certificate grammar.
\end{resultbox}
This task is narrower than the original conjecture and stronger than a numerical packet classification. It asks for an algorithmic cycle formula that exposes exactly how a Hodge packet becomes algebraic.


